\section{The memory block}
\label{sec:method}

\subsection{Architecture}

We take a small GPT ($n_\text{layers}=6$, $d_\text{model}=256$, single head) and
insert one memory module with three interfaces to the residual stream: a
\emph{read} at an early block, a \emph{write} at the last block before the head,
and a per-model \emph{store}. Input text is processed in fixed segments of
$\text{seg\_len}=128$ tokens. Crucially the model does \emph{not} carry a KV cache
across segments: \texttt{forward} is stateless per segment, so the store is the
\emph{only} channel between segments. This makes the controls decisive---with the
store disabled or its contents scrambled, cross-segment recall must fall to
chance, and it does.

\paragraph{Key, query, value.} At layer $\ell_\text{key}=0$ we capture, before the
read, the pre-read hidden state as the key source. With \texttt{key\_pos\_free}
the key uses the token embedding \emph{without} the position embedding, so a key is
a pure function of the token. The read head's query projection is initialised to
the identity and frozen, so at every step query and key are the same function of
the same state and an identical token matches itself exactly. This is the
load-bearing choice: a learned or contextualised key drifts away from the token
and matching collapses to chance (\S\ref{sec:placement}).

\paragraph{Reading.} Retrieval is a batched, annealed Hopfield settling over the
store's keys: cosine scores $q\!\cdot\!K$ are sharpened by a softmax and used to
blend the stored values, which are projected back and injected into the residual
stream through a ReZero gate ($\alpha$ initialised to $0$, so the block is an
exact no-op at initialisation and the store-off baseline is a true control). A
surrogate-gradient trick takes the forward value from a sharp softmax and the
gradient from a soft one, because sharp retrieval and gradient flow are in direct
conflict: on a small store the query gradient vanishes by the second settling step
under a hard schedule.

\paragraph{Writing.} We write one slot per token position (\texttt{write\_per\_token}),
so the store holds exact content rather than a pooled topic vector. A value shift
of $1$ makes the slot keyed on token $j$ return token $j\!+\!1$, i.e.\ the answer
that follows a key. The value is taken at the last block; because the last block's
representation is specialised for next-token prediction, a bare last-block write
degrades, and we add a \emph{normalised block-0 value skip}: the stored value is
the normalised sum of the last-block state and the block-0 state, which restores
the early content the head has abstracted away. This normalised skip is what lets
the objective's ``write at the last block'' work at all, and to scale as the store
grows.

\paragraph{Per-layer read/write improves on read-early/write-late.} The placement
above uses one write (a late block, or the mean of all blocks) and an early read. A
stronger arrangement keeps the \emph{one shared store} but lets \emph{every} layer
both read from it and write its own post-block state into it, so a token contributes
$L$ slots---one per depth---rather than one. Keyed by the same exact position-free
token, this lifts cold single-fact recall from $\recall{80.8}$ (write-the-mean, one
slot per token) to \recall{95.5} and \recall{98.9} over two seeds, and reaches its
recall transition earlier in training. A matched control isolates the cause: keying
each layer's slot by that layer's \emph{evolving residual} instead of the token---the
identical $L\times$ slot count---collapses to \recall{0.2} (the read gates never open,
having nothing to match). So the gain is per-depth \emph{values} that remain
\emph{addressable} because the key is still the exact token, not the extra capacity;
and the block-0 value skip above is subsumed, since layer $0$'s own slot already
carries the early, identity-bearing content. We take per-layer read/write as the
default.

\subsection{Store and persistence}

The store is a fixed-capacity ring buffer, one per example, holding
$[\,\text{keys},\text{values}\,]$ of shape $[B,M,d_\text{mem}]$ plus a running key
mean for centring. It is a non-persistent buffer, so a checkpoint trained at one
capacity loads and runs at another unchanged---store size is a deployment knob,
not a trained constant. A \texttt{save}/\texttt{load} path stamps the file with a
model fingerprint (a foreign model's store is refused, not silently degraded), so
the store survives a process death: we write a store, kill the process, reload it
in a fresh interpreter, and answer from an empty prompt at no loss
(\S\ref{sec:language}).
